% ===================================================================
%  TABELO N° 15 — La Merkato. --- La Manjebli.
%  Feuillets 98 a 102 du fac-simile = folios 96 a 100.
%  Correspondance : folio imprime = numero de feuillet - 2.
%
%  Le feuillet 99 (folio 97) porte en pied de page la signature de
%  cahier « 7. Ido. », qui n'appartient pas au texte et n'est pas
%  transcrite (cf. tabelo 5).
% ===================================================================

% --- Feuillet 98 (folio 96, ouverture de tableau) -------------------
\begin{VUpage}[98]{96}
%%K t15-apar-1 apar
\VUsaut{4.0mm}
\VUtitre{15.0pt}{60}{TABELO N\textsuperscript{o} 15}
\VUsaut{3.0mm}

%%K t15-tit-1 sub
\VUpk{11.4pt}{40}{La Merkato. --- La Manjebli.}
\VUsaut{3.0mm}

%%K t15-01-1 p
\VUblancAlinea
1. --- La maxim multa komercisti qui livras\nl
a ni la vari necesa por nia nutrado su asemblas\nl
sub la \VUgras{halo} \textsuperscript{(1)} dil \VUgras{kovrita merkato} od en la\nl
vicina stradi.

%%K t15-02-1 p
\VUblancAlinea
2. --- La \VUgras{porko-karnisto} \textsuperscript{(2)}, qua vendas\nl
\VUgras{shinko} \textsuperscript{(3)}, \VUgras{sango-sociso} \textsuperscript{(4)}, \VUgras{sociseto} \textsuperscript{(5)}, \VUgras{trip-}\nl
\VUgras{sociseto} \textsuperscript{(6)} e \VUgras{lardo} \textsuperscript{(7)}, stacas proxim la \VUgras{butro-e}\nl
\VUgras{fromajo-vendistino} \textsuperscript{(8)}, di qua l'estaleyo esas\nl
provizita per \VUgras{Holandana fromaji} \textsuperscript{(9)} kun reda\nl
krusto, per \VUgras{Camembert-fromaji} \textsuperscript{(10)}, per\nl
\VUgras{Gruyère-diski} \textsuperscript{(11)} e per fresha \VUgras{butro-bloki} \textsuperscript{(12)}.

%%K t15-03-1 p
\VUblancAlinea
3. --- Avan elca, \VUgras{legum-vendistino} \textsuperscript{(13)} pro\cc
pozas a sua klientini bela \VUgras{tomati} \textsuperscript{(14)}, \VUgras{flor-}\nl
\VUgras{kauli} \textsuperscript{(15)}, \VUgras{salado} \textsuperscript{(16)}, \VUgras{kap-kauli} \textsuperscript{(17)}, \VUgras{kukurbi}\cc
\VUgras{ti} \textsuperscript{(19)}, \VUgras{meloni} \textsuperscript{(18)} e mem \VUgras{fungi} \textsuperscript{(20)}. Ma \VUgras{bir-gar}\cc
\VUgras{sono}, qua haltigis sua \VUgras{livro-veturo} \textsuperscript{(21)} kargita\nl
per \VUgras{bir-bareli} \textsuperscript{(22)}, precize avan elua estaleyo,\nl
impedas elua klienti proximeskar. Gardenis\cc
tino kunportas ad elu korbedo de \VUgras{artichoki} \textsuperscript{(23)}.

%%K t15-tit-2 sub
\VUsaut{4.0mm}
\VUpk{10.2pt}{40}{Vendado e Komprado.}
\VUsaut{3.0mm}

%%K t15-04-1 p
\VUblancAlinea
4. --- En la merkato, l'animo esas granda,\nl
nam arivis la kloko dil \VUgras{auciono-vendo} \textsuperscript{(24)}. La\nl
\VUgras{domo-mastrini} \textsuperscript{(26)}, qui venas adibe por sua\nl
kompri, haltas pasante avan la butiko di la\nl
\VUgras{trip-vendistino} \textsuperscript{(25)}, por komprar \VUgras{stomak-tripi}\nl
o \VUgras{bovyun-kapo.}
\end{VUpage}

% --- Feuillet 99 (folio 97) ------------------------------------------
\begin{VUpage}[99]{97}
%%K t15-05-1 p
\VUblancAlinea
5. --- Grasoza \VUgras{koquistulo} \textsuperscript{(27)} marchandas tre\nl
bel \VUgras{atuno} che la \VUgras{fish-vendistino} \textsuperscript{(28)} qua segun\nl
sua arbitrio plu altigas sua (vendo-)preci. El\nl
recevis granda quanto de \VUgras{anguili, solei, truti}\nl
e \VUgras{karpi.} Elua \VUgras{moruo} e \VUgras{musli} esas tre fresha.

%%K t15-tit-3 sub
\VUsaut{4.0mm}
\VUpk{10.2pt}{40}{Chipa e Chera vari.}
\VUsaut{3.0mm}

%%K t15-06-1 p
\VUblancAlinea
6. --- La \VUgras{pultro-vendistino} \textsuperscript{(29)}, instalita\nl
proxim elu, esas tre konciencoza. Kande el\nl
vendas \VUgras{dindo, ganso, anado} o simpla \VUgras{hanyuno,}\nl
el postulas nur la justa preco.

%%K t15-07-1 p
\VUblancAlinea
7. --- Ye l'angulo di la placo, sur l'unesma\nl
etajo, de poka tempo, esas instalita \VUgras{tel-ven}\cc
\VUgras{disto} \textsuperscript{(30)} che qua la komprantini esas sempre\nl
certa trovar tre bela sortaro de \VUgras{tablo-tuki,}\nl
\VUgras{bok-tuki, avantali} e \VUgras{vishili.} Sur la sama etajo,\nl
\VUgras{kultelisto} \textsuperscript{(31)} instalis sua laboreyo. Il akutigas\nl
la \VUgras{kulteli} di la vendistini en la halo e fabrikas\nl
por la gardenisti \VUgras{serpi} ek la maxim harda\nl
\VUgras{stalo.}

%%K t15-tit-4 sub
\VUsaut{4.0mm}
\VUpk{10.2pt}{40}{La Spicaro.}
\VUsaut{3.0mm}

%%K t15-08-1 p
\VUblancAlinea
8. --- Sub ilua laboreyo, de poka tempo,\nl
esas apertita granda magazino di nutrivi. Ja lu\nl
ne plus esas la simpla e neimportanta \VUgras{spic-}\nl
\VUgras{venderio} \textsuperscript{(32)} olima, qua nur vendis la komu\cc
nuza vari, \VUgras{teo} \textsuperscript{(33)}, \VUgras{chokolado} \textsuperscript{(34)} \VUgras{bloka sukro} \textsuperscript{(37)}\nl
e \VUgras{sapono} \textsuperscript{(38)}. Ibe on vendas irga varo, \VUgras{kra}\cc
\VUgras{kanta kuki}, \VUgras{konservaji} \textsuperscript{(35)}, \VUgras{liquori} \textsuperscript{(36)}, \VUgras{rumo,}\nl
\VUgras{konyako, kirsho, aniz-liquoro} e \VUgras{kuracao.} Ibe\nl
on trovas vildo : \VUgras{leporo} \textsuperscript{(39)}, \VUgras{turdi} \textsuperscript{(40)}, \VUgras{perdri}\cc
\VUgras{ki} \textsuperscript{(41)}, \VUgras{qualii} \textsuperscript{(42)}, \VUgras{sovaj anado} \textsuperscript{(43)} o \VUgras{fazano} \textsuperscript{(44)},\nl
tam facile kam exotika frukti tal quala \VUgras{bana}\cc
\VUgras{ni} \textsuperscript{(46)} ed \VUgras{ananasi.}
\end{VUpage}

% --- Feuillet 100 (folio 98) -----------------------------------------
\begin{VUpage}[100]{98}
%%K t15-09-1 p
\VUblancAlinea
9. --- Spicista \VUgras{garsono} \textsuperscript{(45)} tre komplezema,\nl
esas pronta por donar to quon on deziras :\nl
\VUgras{konfitajo} \textsuperscript{(47)} riba o quinga, \VUgras{graso} \textsuperscript{(48)}, \VUgras{kukom}\cc
\VUgras{breti} \textsuperscript{(49)} o salizita \VUgras{sardini} \textsuperscript{(50)}.

%%K t15-09-2 p
\VUblancAlinea
To esas tre oportuna por la \VUgras{klientini} \textsuperscript{(51)}, qui\nl
trovas, apud la maxim komunuza vari, la pre\cc
mici maxim rara e la frukti maxim saporoza :\nl
sukoza \VUgras{piri} \textsuperscript{(52)}, \VUgras{pruni} \textsuperscript{(53)}, \VUgras{vitberi} \textsuperscript{(54)}, \VUgras{persi}\cc
\VUgras{ki} \textsuperscript{(55)}, \VUgras{mandeli} \textsuperscript{(56)} e \VUgras{nuci} \textsuperscript{(57)}.

%%K t15-tit-5 sub
\VUsaut{4.0mm}
\VUpk{10.2pt}{40}{La Klientaro.}
\VUsaut{3.0mm}

%%K t15-10-1 p
\VUblancAlinea
10. --- La \VUgras{rizo} \textsuperscript{(58)} e la \VUgras{lensi} \textsuperscript{(59)} esas apud\nl
la kesto dil fumizita \VUgras{haringi} \textsuperscript{(60)}; la \VUgras{terpomo} \textsuperscript{(61)}\nl
e la \VUgras{fazeolo} \textsuperscript{(63)} esas vicina al \VUgras{pomi} \textsuperscript{(62)}, \VUgras{figi} \textsuperscript{(64)},\nl
\VUgras{sika pruni} \textsuperscript{(65)} ed \VUgras{avelani} \textsuperscript{(66)}. On trovas en ta\nl
magazino fresha \VUgras{ovi} \textsuperscript{(67)} ed \VUgras{olivi} \textsuperscript{(68)}; pro to\nl
\VUgras{korbi} \textsuperscript{(70)} e \VUgras{provizo-reti} \textsuperscript{(73)} esas rapidege ple\cc
nigita.

%%K t15-11-1 p
\VUblancAlinea
11. --- La \VUgras{kafeo} \textsuperscript{(72)} di ta bonega firmo aqui\cc
ris, inter la \VUgras{servistini} \textsuperscript{(69)} e la \VUgras{koquistini}, justa\nl
reputo, nam la olda \VUgras{spicisto} \textsuperscript{(71)} ipsa maxim\nl
sorgoze torefaktas olu.

%%K t15-tit-6 sub
\VUsaut{4.0mm}
\VUpk{10.2pt}{40}{La Spektaji.}
\VUsaut{3.0mm}

%%K t15-12-1 p
\VUblancAlinea
12. --- Ne omni adiras la merkato por pro\cc
vizar su. En la pitoreska cirkulado sur la stra\cc
deto qua duktas a la halo, on vidas la maxim\nl
diversa personi. Apud afabla \VUgras{bucho-servis}\cc
\VUgras{to} \textsuperscript{(75)}, qua pulsas avan su \VUgras{brueto} \textsuperscript{(74)}, yen du\nl
permisaria soldati tote fiera vidigar sua brilant\nl
uniformo : la \VUgras{husaro} \textsuperscript{(76)} kun blua \VUgras{dolmano} e\nl
\VUgras{plumoza chako}, e la \VUgras{zuavo-kaporalo,} kun bra\cc
cho bufanta e kapkovrita per \VUgras{fezo.}
\end{VUpage}

% --- Feuillet 101 (folio 99) ------------------------------------------
\begin{VUpage}[101]{99}
%%K t15-13-1 p
\VUblancAlinea
13. --- La \VUgras{panifisto} \textsuperscript{(80)}, qua stacas an la solio\nl
di la \VUgras{paniferio} \textsuperscript{(78)} jus petrisis la pasto di sua\nl
\VUgras{pano} \textsuperscript{(79)} e retroprenis del forno la \VUgras{krecenti} \textsuperscript{(81)},\nl
la \VUgras{semli} \textsuperscript{(82)} e la \VUgras{disko-pani} \textsuperscript{(83)} qui esas en la\nl
vetrino.

%%K t15-14-1 p
\VUblancAlinea
14. --- Super la paniferio lojas habila \VUgras{lin}\cc
\VUgras{jistino} \textsuperscript{(84)} qua employas plura \VUgras{brodistini} \textsuperscript{(85)}.\nl
Apud elu habitas \VUgras{lav-glatistino} \textsuperscript{(86)}, qua ami\cc
lizas la \VUgras{linjaro} \textsuperscript{(88)} pos lavir e sikigir ol, e qua\nl
glatigas lu per \VUgras{glatigilo} \textsuperscript{(87)}.

%%K t15-tit-7 sub
\VUsaut{4.0mm}
\VUpk{10.2pt}{40}{Karno, fishi e legumi.}
\VUsaut{3.0mm}

%%K t15-15-1 p
\VUblancAlinea
15. --- En la \VUgras{karnerio} \textsuperscript{(89)}, situata sur la ter\cc
etajo, on vendas omnasorta \VUgras{karni, muton-}\nl
\VUgras{karno} \textsuperscript{(90)}, \VUgras{bov-karno} \textsuperscript{(94)} o \VUgras{bovyun-karno} \textsuperscript{(92)}\nl
kom \VUgras{muton-kruri, bifsteki, rostaji} e \VUgras{kotleti.}\nl
La \VUgras{bucho-garsono} \textsuperscript{(93)} dissekas ta omna karni\nl
ed aparte pozas la \VUgras{medul-osti} \textsuperscript{(96)}. Ye la pordo\nl
di la karnerio esas \VUgras{ostro-vendistino} \textsuperscript{(97)} qua\nl
vendas \VUgras{ostri} \textsuperscript{(98)}. El apertas oli se on deziras\nl
lo.

%%K t15-16-1 p
\VUblancAlinea
16. --- Ta omna komercisti pagas patento o\nl
plaso-taxo; pro to la \VUgras{policisto} \textsuperscript{(100)} senkompate\nl
persekutas la kompatinda \VUgras{legum-kolportisti}\cc
\VUgras{ni} \textsuperscript{(99)} qui konkurencas li disportante per sua\nl
vetureti \VUgras{karoti, rafaneti, napi, poreli} od \VUgras{aspa}\cc
\VUgras{raga fasketi, fresha pizi, spinati, oxalo, kreso}\nl
e \VUgras{petroselo.}

%%K t15-tit-8 sub
\VUsaut{4.0mm}
\VUpk{10.2pt}{40}{Atencez la policisto !}
\VUsaut{3.0mm}

%%K t15-17-1 p
\VUblancAlinea
17. --- La puero qua vendas \VUgras{alio} \textsuperscript{(101)} ed\nl
\VUgras{onyoni,} tre bone savas, ke anke lu, ne darfas
\parplein
\end{VUpage}

% --- Feuillet 102 (folio 100) -----------------------------------------
\begin{VUpage}[102]{100}
%%K t15-17-1 p suite
\VUcontinue
vendar sua vari proxim la merkato. Il fugas,\nl
kurante kun risko renversor la korbo de \VUgras{krabi,}\nl
\VUgras{kankri} e \VUgras{kreveti} di la \VUgras{vendisto} \textsuperscript{(102)} di \VUgras{langusti}\nl
e \VUgras{homardi.}

%%K t15-18-1 p
\VUblancAlinea
18. --- Omni agitas su e bruisegas; ma to\nl
ne impedas audar klare la voco dil kolportanta\nl
\VUgras{vitristo} \textsuperscript{(103)}, nek la klamo dil \VUgras{karbonisto} \textsuperscript{(104)}\nl
qua jus livis por kelka instanti sua chareto por\nl
livrar \VUgras{karbona sakedo} \textsuperscript{(105)}.

\VUsaut{6.0mm}
\VUfilet{22mm}
\end{VUpage}
